\chapter*{}
\thispagestyle{empty}
\cleardoublepage

\thispagestyle{empty}

\input{portada/portada_2}



\cleardoublepage
\thispagestyle{empty}

\begin{center}
{\large\bfseries ReMind: Plataforma web para estimulación cognitiva de personas con Alzheimer}\\
\end{center}
\begin{center}
Alberto Siles Molina\\
\end{center}

\vspace{0.7cm}
\noindent{\textbf{Palabras clave}: Alzheimer, estimulación cognitiva, aplicación web, Angular, Spring Boot, minijuegos, terapia no farmacológica}\\

\vspace{0.7cm}
\noindent{\textbf{Resumen}}\\
Este Trabajo Fin de Grado presenta ReMind, una aplicación web diseñada para ofrecer estimulación cognitiva a personas mayores, especialmente aquellas en fases tempranas de Alzheimer o en riesgo de desarrollarlo. La plataforma integra seis minijuegos terapéuticos basados en ejercicios validados científicamente, que trabajan áreas como la memoria, la atención, el lenguaje, el cálculo y las funciones ejecutivas. El sistema sigue una arquitectura cliente-servidor de tres capas, con Angular en el frontend y Spring Boot en el backend, junto con una base de datos MySQL. Además de los juegos, la aplicación incluye un sistema de gestión de usuarios con tres roles (administrador, terapeuta y paciente), permitiendo a los terapeutas asignar agendas personalizadas y realizar un seguimiento del progreso de cada paciente. El diseño de la interfaz sigue los estándares WCAG 2.1 para garantizar la accesibilidad del público objetivo.\\

\cleardoublepage


\thispagestyle{empty}


\begin{center}
{\large\bfseries ReMind: Web Platform for Cognitive Stimulation of People with Alzheimer's Disease}\\
\end{center}
\begin{center}
Alberto Siles Molina\\
\end{center}

\vspace{0.7cm}
\noindent{\textbf{Keywords}: Alzheimer, cognitive stimulation, web application, Angular, Spring Boot, minigames, non-pharmacological therapy}\\

\vspace{0.7cm}
\noindent{\textbf{Abstract}}\\
This Bachelor's Thesis presents ReMind, a web application designed to provide cognitive stimulation for older adults, particularly those in early stages of Alzheimer's disease or at risk of developing it. The platform includes six therapeutic minigames based on scientifically validated exercises, targeting cognitive areas such as memory, attention, language, arithmetic and executive functions. The system follows a three-layer client-server architecture, using Angular for the frontend and Spring Boot for the backend, with a MySQL database. In addition to the games, the application includes a user management system with three roles (administrator, therapist and patient), allowing therapists to assign personalised schedules and monitor each patient's progress. The interface design follows WCAG 2.1 standards to ensure accessibility for the target audience.\\

\chapter*{}
\thispagestyle{empty}

\noindent\rule[-1ex]{\textwidth}{2pt}\\[4.5ex]

Yo, \textbf{Alberto Siles Molina}, alumno de la titulación Ingeniería Informática de la \textbf{Escuela Técnica Superior
de Ingenierías Informática y de Telecomunicación de la Universidad de Granada}, con DNI 26050580K, autorizo la
ubicación de la siguiente copia de mi Trabajo Fin de Grado en la biblioteca del centro para que pueda ser
consultada por las personas que lo deseen.

\vspace{6cm}

\noindent Fdo: Alberto Siles Molina

\vspace{2cm}

\begin{flushright}
Granada a 22 de junio de 2026 .
\end{flushright}


\chapter*{}
\thispagestyle{empty}

\noindent\rule[-1ex]{\textwidth}{2pt}\\[4.5ex]

D. \textbf{Nuria Medina Medina}, Profesor del Área de  del Departamento de Lenguajes y Sistemas Informáticos de la Universidad de Granada.

\vspace{0.5cm}

\textbf{Informan:}

\vspace{0.5cm}

Que el presente trabajo, titulado \textit{\textbf{ReMind: Plataforma web para estimulación cognitiva de personas con Alzheimer}},
ha sido realizado bajo su supervisión por \textbf{Alberto Siles Molina}, y autorizamos la defensa de dicho trabajo ante el tribunal
que corresponda.

\vspace{0.5cm}

Y para que conste, expiden y firman el presente informe en Granada a 22 de junio de 2026 .

\vspace{1cm}

\textbf{Los directores:}

\vspace{5cm}

\noindent \textbf{Nuria Medina Medina}

\chapter*{Agradecimientos}
\thispagestyle{empty}

Quiero dedicar este trabajo a mis padres, por su apoyo incondicional durante todos estos años, por confiar en mí incluso cuando yo mismo dudaba y por enseñarme que el esfuerzo siempre tiene recompensa. A mi hermano, por estar siempre ahí, por los momentos de desconexión y por recordarme que hay vida más allá de la pantalla.

A los amigos que he hecho durante la carrera, por hacer de esta etapa algo mucho más llevadero. Por las sesiones de estudio interminables, los trabajos en grupo, los cafés, las risas y también las quejas compartidas. Sin ellos, este camino habría sido mucho más difícil y mucho menos divertido.

También quiero agradecer a mi tutor por la orientación y los consejos durante el desarrollo de este proyecto, por estar disponible cuando lo he necesitado y por ayudarme a dar forma a esta idea.

Gracias a todos por formar parte de esto.

